%!TEX program = lualatex
\documentclass[ngerman,headheight=70pt]{scrartcl}
\usepackage[ngerman]{babel}
\usepackage{blindtext}
\usepackage[utf8]{luainputenc}
\usepackage{scrlayer-scrpage}
\usepackage[T1]{fontenc}
\usepackage{fontspec}
\usepackage{multicol}
\usepackage{enumitem}
\usepackage{epigraph}
\usepackage[german=quotes]{csquotes}
\usepackage[bottom=2cm,footskip=8mm]{geometry}
\setmainfont{Source Sans Pro}
\MakeOuterQuote{"}

\setlength{\parindent}{0pt}
\setlength{\parskip}{6pt}

\pagestyle{plain.scrheadings}
\thispagestyle{scrheadings}
\lohead{\vspace*{-1.5cm}\includegraphics[scale=0.4]{../images/stupa_logo_rot.pdf}}
\rohead{\\Präsidium des\\Studierendenparlaments}

\lofoot[Universität Hamburg · Präsidium des Studierendenparlaments\\\hrule Von-Melle-Park 5 · D-20146 Hamburg]{Universität Hamburg · Präsidium des Studierendenparlaments\\\hrule Von-Melle-Park 5 · D-20146 Hamburg}
\cofoot[]{}
\cohead[\thepage]{}

\renewcommand{\thesubsection}{TOP \arabic{subsection}}
\renewcommand{\thesubsubsection}{\arabic{subsubsection}.}
\setcounter{subsection}{-1}

\begin{document}
    UHH · StuPa-Präsidium · Von-Melle-Park 5 · D-20146 Hamburg

    \section*{Beschluss des Studierendenparlamentes vom 16. Juni 2016}
    \subsection*{Antrag auf Beitritt zum FZS}

    Der "freie zusammenschluss von studentInnenschaften" (fzs) ist als
    Dachverband der Studierendenschaften die institutionelle Plattform, über die
    sich eine kritische und eingreifende Studierendenschaft in der Bundesrepublik
    organisieren kann. Sowohl über die regelmäßigen Mitgliederversammlungen, als
    auch in inhaltlichen Ausschüssen und Seminaren gibt es Orte, über die Austausch
    stattfinden und verbindliche Zusammenarbeit sich entwickeln kann. Dabei kann
    es, wie beispielsweise im Bereich der Sozialpolitik, um spezifisch-fachliche
    Belange gehen oder sich allgemein mit der Frage nach den Möglichkeiten
    studentischer Politik befassen.

    Als Verfasste Studierendenschaft der Uni Hamburg sehen wir uns in der
    Verantwortung die erreichten Auseinandersetzungen und die erkämpften
    Verbesserungen in Hamburg für das Bundesgebiet zu verallgemeinern. Exemplarisch
    steht dafür die Studienreform für Bildung durch Wissenschaft und zur
    Abschaffung des Bachelor-Master-Systems. Als Verfasste Studierendenschaft
    haben wir uns an der Uni Hamburg einer Studienreform verschrieben, in der
    Studium im Allgemeinen die Entfaltung der eigenen Persönlichkeit durch die
    (exemplarische) kritische Auseinandersetzung mit den gesellschaftlichen
    Verhältnissen bedeutet. Damit treten wir in den Konflikt mit der Realität an
    den Hochschulen: So ist die sog. Bologna-Reform mit der in der jüngeren
    Vergangenheit so wirkmächtigen, aber heftig gescheiterten, neoliberalen
    Erzählung, dass das Ende der Geschichte erreicht sei und jede*r sich für das
    eigene Glück (freilich im Beruf) nur redlich anzustrengen brauche, versucht
    worden zu legitimieren. Dagegen gilt es, sich zusammen zu organisieren,
    gemeinsame Handlungsfähigkeit zu gewinnen und das Studium wieder in den Dienst
    der eigentlichen Aufgabe zu stellen, nämlich gesellschaftliche Verbesserungen
    zu initiieren und zu flankieren.

    Diese Politik, die auf die Bildung mündiger Menschen zielt muss darauf abzielen,
    dass jede*r die Möglichkeit wahrnehmen kann, sich für die solidarische
    Wahrnehmung der gemeinsamen Interessen zu organisieren. Dieser Ansatz bringt
    as Erfordernis einer starken Schwerpunktsetzung auf die Auseinandersetzungen
    "vor Ort" mit sich und begrenzt Stellvertreterpolitik. Vor diesem Hintergrund
    beauftragt das Studierendenparlament der Universität Hamburg den AStA,
    aufbauend auf der studentischen Politik an der Uni Hamburg, sich um die
    Mitgliedschaft im fzs bemühen.

    Es ist prinzipiell sinnvoll, sich
    zusammenzuschließen und für das gemeinsame Interesse zu streiten. Gerade vor
    dem Hintergrund einer Wissenschaftspolitik, die zunehmend bundesweit --
    beispielsweise in Gestalt der KMK – organisiert wird, erscheint es notwendig,
    dass auf den entsprechenden Ebenen eine verbindliche und inhaltlich fundierte
    politische Organisation besteht. Der Blick auf die verbindenden Einschränkungen
    -- u.a. Bachelor/Master, die Exzellenzinitiative (jetzt Exzellenzstrategie),
    strukturelle Unterfinanzierung -- und ihre politischen Ursachen offenbart die
    grundsätzliche Notwendigkeit, aus den eigenen regionalen Kontexten
    herauszutreten. Es bedarf einer überregionalen Koordination und der
    Möglichkeiten, von den spezifischen Auseinandersetzungen voneinander zu
    lernen, um gemeinsam schlagfertig zu werden.

    Immer deutlicher wird dabei, dass die Politik der Unternehmerischen Hochschule
    von immer mehr Menschen abgelehnt wird – kritische Studierendenschaften sind
    es an den Hochschulen, die eine wissenschaftspolitische und kulturelle
    Alternative zum Status Quo entwerfen müssen.
    \textit{Dieser Aufgabe wollen wir uns, als Verallgemeinerung unser Stände in
    Hamburg, verstärkt mit den im fzs organisierten Kommiliton*innen stellen!}
\end{document}
